% Section 6.3, from lectures/L06/README.md: the objectives, the evaluation questions, and the next
% lecture.

\needspace{10\baselineskip}
\section{Review}
\label{c6:sec:review}

\subsection*{What you should be able to do now}
After this chapter, you should:
\begin{itemize}
  \item Know why regular neural networks aren't well suited to image classification.
  \item Know at a high level how convolutional layers are structured.
  \item Be able to extract features from images via kernels.
  \item Be able to feed an input through a small CNN by hand (feedforward).
  \item Have started implementing a conv layer in software (\code{ml::ConvLayer}).
\end{itemize}

\needspace{6\baselineskip}
\subsection*{Questions to test yourself}
You should be able to explain:
\begin{enumerate}
  \item Why are regular neural networks (dense layers only) less suited to image classification
    than a CNN?
  \item Can you, in your own words, describe the purpose of each layer type in a CNN:
    convolutional, pooling, flatten, and dense?
  \item Without looking at \secref{c6:sec:cnn}, can you describe how a kernel is applied across an
    image during feedforward in a conv layer?
  \item Which parts of your \code{ml::ConvLayer} so far correspond to feedforward?
\end{enumerate}

\needspace{6\baselineskip}
\subsection*{Looking ahead}
Chapter~\ref{c7:ch:layers} finishes the conv layer, then implements max pooling and flatten layers
in software.
